\documentclass[11pt,oneside,openany]{book}
\usepackage[paperwidth=6.5in,paperheight=9.5in,inner=.72in,outer=.72in,top=.75in,bottom=.78in,headheight=14pt]{geometry}
\usepackage{fontspec}
\setmainfont{lmroman10-regular.otf}[BoldFont=lmroman10-bold.otf,ItalicFont=lmroman10-italic.otf,BoldItalicFont=lmroman10-bolditalic.otf,SmallCapsFont=lmromancaps10-regular.otf]
\usepackage[ngerman]{babel}
\usepackage{amsmath,amssymb,graphicx,microtype,fancyhdr,needspace}
\usepackage{tikz}
\usepackage[unicode,hidelinks]{hyperref}
\hypersetup{pdftitle={Carl Friedrich Gauss: Band III, Vollständige Dissertation von 1799},pdfauthor={Carl Friedrich Gauss}}
\setlength{\parindent}{1.2em}\setlength{\parskip}{0pt plus .3pt}
\setlength{\emergencystretch}{1.4em}
\pagestyle{fancy}\fancyhf{}
\fancyhead[L]{\small Ein neuer Beweis}
\fancyfoot[C]{\thepage}\renewcommand{\headrulewidth}{.25pt}
\newcommand{\SourceLeaf}[1]{\hypertarget{native-#1}{}\label{native-#1}}
\newcommand{\Article}[1]{\par\medskip\Needspace{4\baselineskip}\pdfbookmark[1]{Art. #1}{first-paper-art-#1}\centerline{\textbf{#1.}}\nobreak\smallskip}
\renewcommand{\thefootnote}{*)}
\newcommand{\HistoricalNote}[1]{\footnote{\csname HistoricalNoteText#1\endcsname}}
\widowpenalty=10000\clubpenalty=10000\displaywidowpenalty=10000
\raggedbottom
\newcommand{\BareRoot}[2][]{\mathord{{}^{#1}\!\surd}#2}
